The descriptive differences between the regional reruns and the full-montage reference were next tested for every electrode using session-paired comparisons. For each electrode in the four anatomical subsets shown in Figure~\ref{fig:eeg_map}, its session-level $F_1$ in the rerun of its own region was paired with its session-level $F_1$ in the full-montage run, using the sessions common to both. For example, the $F_1$ of Fp1 when the pipeline was run on the frontal electrodes only was compared with the $F_1$ of Fp1 when all 32 channels were supplied to the pipeline, and the same comparison was made for every other electrode within its own region. Each electrode was tested with a two-tailed Wilcoxon signed-rank test paired by session, with $p$-values Bonferroni-corrected over the electrodes tested within each dataset. Figure~\ref{fig:exp1_subset_per_electrode_stats} shows, for every electrode, its session-mean $F_1$ in the regional rerun overlaid on its session-mean $F_1$ in the full-montage run, with an asterisk marking electrodes that remained significant after correction.

On Internal (Figure~\ref{fig:exp1_subset_per_electrode_stats}A), 11 of the 24 electrodes remained significantly different after correction, and every significant difference was a reduction in the regional rerun. Despite a visible drop in their session means, neither of the two highest-performing electrodes, Fp2 and Fp1, differed significantly from the full-montage run (adjusted $p > 0.05$ for both) In contrast, AF4 and AF3, the third- and fourth-highest-performing electrodes, were significantly lower in the regional rerun. In the central region, C3 showed a significant reduction whereas C4 did not, and the pronounced C3--C4 asymmetry in absolute $F_1$ persisted. Several parietal and occipital electrodes, namely CP2, CP5, CP6, PO3, and PO4, also reached significance despite their very low absolute $F_1$, whereas most of the remaining posterior electrodes did not.

On Cao2018 (Figure~\ref{fig:exp1_subset_per_electrode_stats}B), significant differences were far more widespread. Sixteen of the 18 electrodes remained significantly different after correction, and again every significant difference was a reduction. As on Internal, only the frontopolar pair Fp1 and Fp2 did not reach significance (adjusted $p > 0.05$ for both), whereas F3 and F4, the third- and fourth-highest-performing electrodes, did. Unlike on Internal, every central, parietal, and occipital electrode reached significance, including both C3 and C4.

Taken together, the electrode-level comparisons in Figure~\ref{fig:exp1_subset_per_electrode_stats} show that a visibly lower descriptive mean did not by itself imply a statistically significant change. Fp1 and Fp2, the two highest-performing electrodes in each dataset, showed no significant paired difference after correction in either dataset (adjusted $p > 0.05$), whereas 11 of the other 22 electrodes on Internal and all 16 on Cao2018 were significantly lower in the regional rerun (adjusted $p < 0.05$). Significant differences were therefore more widespread on Cao2018 than on Internal, where 11 of the 18 non-frontal electrodes did not reach significance despite showing the same direction of change.
%These tests probe a configuration effect, whether an electrode's own detection performance changes when its region operates alone rather than as part of the complete montage. That question is distinct from the spatial pattern described above, in which frontal electrodes simply detect blinks better than parietal and occipital electrodes regardless of which subset scored them.


\begin{figure}[ht]
    \centering
    \includegraphics[width=\linewidth]{fig_exp1_subset_per_electrode_stats.pdf}
    \caption{Per-electrode comparison of $F_1$ for each electrode in the rerun of its own anatomical region (narrow, opaque bars) and in the full-montage run (wide, translucent bars), for (A)~Internal and (B)~Cao2018. Bars show session-mean $F_1$. Electrodes are grouped by the anatomical regions in Figure~\ref{fig:eeg_map} and ordered by descending regional-rerun $F_1$ within each region. Comparisons are paired by session using two-tailed Wilcoxon signed-rank tests, with $p$-values Bonferroni-corrected separately over the electrodes tested within each dataset. An asterisk marks adjusted $p < 0.05$.}
    \label{fig:exp1_subset_per_electrode_stats}
\end{figure}